\appendix

\section{Proof Details}

\subsection{Step ID: Step Goal}\label{app:step-id}

Each appendix subsection should contain the full paper-ready translation of the
corresponding step derivation. Public appendix proof content should use only
lemma, theorem, proposition, corollary, and proof environments, apart from
appendix headings and concise mathematical exposition needed to introduce a
result. Convert any step-local claim, hypothesis, or invariant into one of the
allowed public theorem-style units, or integrate it into a proof environment;
do not emit public claim, hypothesis, or invariant environments or proof
dependencies. Do not replace detailed derivations with a summary. When choosing
between more helper notation and more visible derivation, prefer the visible
derivation unless the notation clearly improves repeated complex calculations
without hiding any proof obligation.

\begin{lemma}[Local Result Title]\label{lem:step-id-unit-001}
Under Assumptions~\ref{assump:example} and any named prior theorem-style
results used by this unit, if the local conditions for this unit hold, then
state the exact conclusion in current notation.
\end{lemma}

\begin{proof}
Translate the full local derivation for this unit into LaTeX. Preserve the
definitions, inequalities, recursions, case checks, constants, parameter
dependence, and cited-result assumption discharge needed to audit the argument.
\end{proof}

\begin{proposition}[Step Assembly]\label{prop:step-id-assembly}
Under the numbered assumptions and preceding theorem-style results used by this
step, if the step-local conditions hold, then state the exact step-level
conclusion as a theorem-style result.
\end{proposition}

\begin{proof}
Combine the preceding local units and earlier labeled appendix results to prove
the step-level conclusion. Use internal references such as
Lemma~\ref{lem:step-id-unit-001}; use BibTeX-backed citation commands only for
external sources.
\end{proof}

\begin{proposition}[Rate Specialization Bridge]\label{prop:rate-specialization-bridge}
Use this proposition only when the public main theorem or corollary simplifies a
technical appendix rate. State the auxiliary parameter choices, the public
assumptions, the technical theorem conditions being verified, and the final
simplified rate.
\end{proposition}

\begin{proof}
Verify every technical condition from the source theorem, prove each term
absorption or simplification by an explicit inequality under stated
admissibility thresholds, convert the probability statement to the public
probability mode, and state the final hidden-constant dependence.
\end{proof}

\subsection{Proof of the Main Theorem}

\begin{proof}[Proof of the main theorem]
Combine the earlier theorem-style results and direct algebraic implications to
prove Theorem~\ref{thm:main}. Cite only lemmas, theorems, propositions,
corollaries, or numbered assumptions as proof authorities; do not use public
claims, hypotheses, invariants, subsection titles, or local unit IDs as proof
authorities.
\end{proof}
